\documentclass[11pt]{article}

\usepackage[T1]{fontenc}
\usepackage{lmodern}
\usepackage{microtype}
\usepackage{amsmath,amssymb,amsthm,mathtools}
\usepackage{booktabs}
\usepackage{geometry}
\usepackage{xcolor}
\usepackage{hyperref}
\usepackage[nameinlink,noabbrev]{cleveref}
\usepackage{tikz}
\usepackage{pgfplots}
\usepackage{pgfplotstable}
\usepgfplotslibrary{groupplots}

\geometry{margin=1in}
\pgfplotsset{compat=1.18}
\makeatletter
\renewcommand{\@seccntformat}[1]{\csname the#1\endcsname\hspace{1.25em}}
\makeatother
\hypersetup{
  colorlinks=true,
  linkcolor=blue!55!black,
  citecolor=green!45!black,
  urlcolor=blue!60!black
}

\definecolor{Teal}{HTML}{2A9D8F}
\definecolor{Coral}{HTML}{E76F51}
\definecolor{Amber}{HTML}{E9A23B}
\definecolor{Blue}{HTML}{457B9D}
\definecolor{Dark}{HTML}{264653}
\definecolor{Pale}{HTML}{EAF4F2}

\newtheorem{theorem}{Theorem}[section]
\newtheorem{proposition}[theorem]{Proposition}
\newtheorem{lemma}[theorem]{Lemma}
\theoremstyle{definition}
\newtheorem{definition}[theorem]{Definition}
\newtheorem{remark}[theorem]{Remark}
\renewcommand{\qedsymbol}{}

\newcommand{\Ssurf}{\mathcal S}
\newcommand{\Lfold}{\mathcal L_{\mathrm{fold}}}
\newcommand{\LN}{\mathcal L_N}
\newcommand{\Csym}{\mathcal C_0}
\newcommand{\Crat}{\mathcal C_1}
\newcommand{\DiscF}{\Delta_F}
\newcommand{\cb}{0.613055096856276}
\newcommand{\cabs}{0.673741152200546}

\pgfmathdeclarefunction{ufold}{1}{%
  \pgfmathparse{-(1-#1)*#1*(2-#1)*(1+2*#1-#1*#1)/(#1*#1*#1-3*#1+1)}%
}
\pgfplotstableread{
c Sone Stwo Sthree
0.613055096856 0.832749933396 0.000000000000 0.553649300937
0.622728719435 0.834556278962 0.019816171686 0.550566196369
0.632402342013 0.836333223457 0.039485267032 0.546797634441
0.642075964592 0.838089816529 0.058969659229 0.542335725100
0.651749587171 0.839835467001 0.078229521361 0.537174580890
0.661423209749 0.841579989780 0.097222980237 0.531310467539
0.671096832328 0.843333644368 0.115906289018 0.524741933185
0.680770454906 0.845107164016 0.134234014798 0.517469912748
0.690444077485 0.846911774681 0.152159236429 0.509497804388
0.700117700064 0.848759203138 0.169633747004 0.500831515547
0.709791322642 0.850661673801 0.186608254608 0.491479476721
0.719464945221 0.852631894040 0.203032574186 0.481452621850
0.729138567799 0.854683028057 0.218855802690 0.470764335077
0.738812190378 0.856828659670 0.234026469039 0.459430364427
0.748485812957 0.859082744598 0.248492649831 0.447468703837
0.758159435535 0.861459553140 0.262202041099 0.434899445789
0.767833058114 0.863973604357 0.275101975677 0.421744607498
0.777506680692 0.866639593077 0.287139374800 0.408027934277
0.787180303271 0.869472311201 0.298260621175 0.393774684193
0.796853925850 0.872486564883 0.308411338850 0.379011398467
0.806527548428 0.875697089234 0.317536062266 0.363765662300
0.816201171007 0.879118462170 0.325577772592 0.348065860816
0.825874793585 0.882765018984 0.332477273079 0.331940934721
0.835548416164 0.886650769107 0.338172365693 0.315420140009
0.845222038743 0.890789316390 0.342596777231 0.298532815677
0.854895661321 0.895193784052 0.345678761940 0.281308162942
0.864569283900 0.899876745235 0.347339275387 0.263775038924
0.874242906478 0.904850159896 0.347489564216 0.245961767145
0.883916529057 0.910125318562 0.346027936592 0.227895966638
0.893590151635 0.915712793224 0.342835347473 0.209604400832
0.903263774214 0.921622395485 0.337769211020 0.191112846842
0.912937396793 0.927863141863 0.330654460347 0.172445985239
0.922611019371 0.934443226009 0.321270146449 0.153627309963
0.932284641950 0.941369997445 0.309328432834 0.134679057583
0.941958264528 0.948649946330 0.294439801385 0.115622154830
0.951631887107 0.956288693684 0.276051227924 0.096476183040
0.961305509686 0.964290986410 0.253325650373 0.077259357957
0.970979132264 0.972660696471 0.224874410979 0.057988523232
0.980652754843 0.981400823513 0.188032833595 0.038679155865
0.990326377421 0.990513500257 0.136046911071 0.019345381807
1.000000000000 1.000000000000 0.000000000000 0.000000000000
}\StokesTable
\pgfplotstableread{
t u N M
0.910874496227 1.000000000000 -0.624163448219 -0.261919663769
0.898570984577 0.959139710045 -0.581378114160 -0.215643645946
0.886145789889 0.920223038547 -0.541415629433 -0.171349104722
0.873593251028 0.883053718163 -0.504022685451 -0.128830295828
0.860907328993 0.847460468518 -0.468978566000 -0.087906987128
0.848081569503 0.813293131759 -0.436090113271 -0.048420510428
0.835109060733 0.780419500131 -0.405187597707 -0.010230519456
0.821982385398 0.748722694997 -0.376121308190 0.026787689497
0.808693566249 0.718098988408 -0.348758720463 0.062745405459
0.795234003814 0.688455982209 -0.322982132885 0.097742365972
0.781594405009 0.659711077836 -0.298686682260 0.131868323878
0.767764700936 0.631790183808 -0.275778670687 0.165204371800
0.753733951759 0.604626618697 -0.254174148320 0.197824067463
0.739490236135 0.578160175658 -0.233797707886 0.229794394496
0.725020521980 0.552336321167 -0.214581455302 0.261176586678
0.710310514593 0.527105505732 -0.196464127457 0.292026838373
0.695344477089 0.502422568435 -0.179390333591 0.322396919696
0.680105016704 0.478246220419 -0.163309900918 0.352334711682
0.664572828708 0.454538595041 -0.148177308581 0.381884674029
0.648726387174 0.431264854489 -0.133951196774 0.411088255886
0.632541568472 0.408392844410 -0.120593940108 0.439984258395
0.615991188697 0.385892789426 -0.108071276087 0.468609156320
0.599044429712 0.363737023590 -0.096351981118 0.496997384910
0.581666119164 0.341899750734 -0.085407587616 0.525181597268
0.563815816366 0.320356830426 -0.075212136885 0.553192896649
0.545446635986 0.299085585895 -0.065741963196 0.581061047567
0.526503711372 0.278064630785 -0.056975505275 0.608814668997
0.506922152705 0.257273712046 -0.048893141962 0.636481412593
0.486624280980 0.236693566612 -0.041477049298 0.664088128443
0.465515796881 0.216305789820 -0.034711076761 0.691661020652
0.443480336093 0.196092713769 -0.028580640689 0.719225794772
0.420371493889 0.176037294017 -0.023072633280 0.746807798937
0.396000713888 0.156123003167 -0.018175345813 0.774432160429
0.370118073383 0.136333730056 -0.013878404974 0.802123919273
0.342380096510 0.116653683323 -0.010172721402 0.829908160423
0.312291959535 0.097067298252 -0.007050449735 0.857810146022
0.279093681921 0.077559145819 -0.004504959653 0.885855449245
0.241505000911 0.058113842907 -0.002530817533 0.914070091240
0.197027688516 0.038715962701 -0.001123778549 0.942480682709
0.139206748637 0.019349944226 -0.000280789141 0.971114571794
0.000000000000 0.000000000000 0.000000000000 1.000000000000
}\CriticalTable

\title{\textbf{Observation Discriminant and Spherical Fold Image}\\
\large of the Orthogonal-Circle Ruled Surface}
\author{Akari Hayami (Jian-Yu Huang)}
\date{July 2026}

\begin{document}
\maketitle

\begin{abstract}
We determine the full critical locus of the planar observation map attached
to the orthogonal-circle ruled surface.  Its Jacobian factors as
\[
J_F=-\frac{2\sin t}{Q^3}\bigl(A(c)+uB(c)\bigr),
\]
so the critical set is not only the previously isolated rational branch: it
also contains the complete symmetry segment $t=0$.  The symmetry segment is
an ordinary Whitney fold for $u>0$.  The two signed rational arms are ordinary
folds away from their common initial point and occupy
$c_b\leq c\leq1$, where
$c_b=0.613055096856276\ldots$ is the unique root in $(0,1)$ of
$c^5-5c^4+6c^3-c^2+c-1$.  At $(t,u)=(0,0)$ the two critical branches meet and
the ordinary-fold condition fails; locally
$J_F=2t(u-t^2)+\text{higher-order terms}$.  Their critical values are
different objects from the normalized surface-position curve: the symmetry
fold maps to the segment $N=0$, while the rational arms share one
discriminant curve tangent to that segment at $(N,M)=(0,1)$.

For completeness we retain the parameter-free spherical radial image
$\Phi=S/\lVert S\rVert$ of the rational branch.  It consists of two mirrored
arcs that start regularly at $(1,0,0)$ and meet the parameter boundary at
\[
\left(\sqrt{\frac{2-c_b}{2}},0,\sqrt{\frac{c_b}{2}}\right)
\approx(0.8327499334,0,0.5536493009).
\]
This radial image is not identified with the observation discriminant or the
Gauss map.  Its third spherical component is monotone and maximal at the
boundary.  Exact polynomial identities, Sturm counts and numerical checks are
reproduced by a standard-library verifier.  No Jones-spinor, quantum,
spinorial or Atiyah--Patodi--Singer interpretation is claimed.
\end{abstract}

\section{Correction scope and claim status}

Let $c=\cos t$, $s=\sin t$ and
$Q=\sqrt{c(2-c)}$, with
$t\in[-\pi/2,\pi/2]$ and $u\in[0,1]$.  The ruled surface of
\cite{hayami2026surface} is
\begin{equation}\label{eq:surface}
S(t,u)=
\begin{pmatrix}
c+2u(1-c)\\
(1-u)s\\
uQ
\end{pmatrix}.
\end{equation}
Four distinct objects must be kept separate:
\begin{enumerate}
  \item the component-zero curve
        $\LN=\{N=0\}$, given away from $c=1$ by
        $u=(1-c)/2$;
  \item the isolated immersion zeros, obtained only after all components of
        $S_t\times S_u$ vanish; and
  \item the full critical locus $\operatorname{Crit}(F)=\{J_F=0\}$ of the
        planar observation map;
  \item the discriminant $F(\operatorname{Crit}(F))\subset\mathbb R^2$ and
        the separate radial sphere map $S/\lVert S\rVert\subset S^2$.
\end{enumerate}
Only the image of the third object under $F$ is the observation
discriminant.  Radial normalization of the surface position is a different,
origin-dependent construction.

\begin{center}
\begin{tabular}{@{}p{0.37\textwidth}p{0.18\textwidth}p{0.34\textwidth}@{}}
\toprule
Statement & Status here & Input\\
\midrule
Full critical-locus decomposition & theorem & exact factorization of $J_F$\\
Fold status away from $(0,0)$ & theorem & kernel transversality/Sturm count\\
Degenerate intersection at $(0,0)$ & theorem & local expansion\\
Observation discriminant & theorem & restriction of $F$\\
Spherical radial fold image & theorem & \eqref{eq:surface}\\
\addlinespace[2pt]
Jones, spinor or APS interpretation & not claimed & absent from this geometry\\
\bottomrule
\end{tabular}
\end{center}

\section{Observation map and full critical locus}

The complex observation field used in \cite{hayami2026surface} is
$F=N+iM$, where
\begin{align}
N(c,u)&=(1-c)(1-c-2u),\\
M(c,u)&=
\frac{(1-u)c^2(2-c)-u(1-c)^2(1+c)}{Q}.
\end{align}
Its fold Jacobian is
$J_F=N_tM_u-N_uM_t$.

\begin{theorem}[Full critical locus]\label{thm:full-critical}
Where $Q>0$, the Jacobian factors exactly as
\begin{equation}\label{eq:jac-factor}
J_F=-\frac{2s}{Q^3}\bigl(A(c)+uB(c)\bigr),
\end{equation}
where
\[
A(c)=(1-c)c(2-c)(1+2c-c^2),
\qquad
B(c)=c^3-3c+1.
\]
Consequently the critical locus in the parameter strip has two components:
\begin{align}
\Csym&=\{(t,u):t=0,\ 0\leq u\leq1\},\\
\Crat&=\{(t,u):A(\cos t)+uB(\cos t)=0\}.
\end{align}
\end{theorem}

\begin{proof}
Direct differentiation and collection in powers of $u$ gives
\eqref{eq:jac-factor}.  In the real parameter interval, $s=0$ is exactly
$t=0$.  The other factor gives $\Crat$.  Neither factor may be cancelled when
determining the full zero set.
\end{proof}

\begin{proposition}[Rational critical branch]\label{thm:true-fold}
Away from $s=0$ and the vertical asymptote $B(c)=0$, the component
$\Crat$ is
\begin{equation}\label{eq:true-fold}
u=u_{\mathrm{fold}}(c):=
-\frac{(1-c)c(2-c)(1+2c-c^2)}
       {c^3-3c+1}.
\end{equation}
The curve $u=(1-c)/2$ is instead $\LN=\{N=0\}$ and is not equal to
\eqref{eq:true-fold}.
\end{proposition}

\begin{proof}
Solving $A+uB=0$ proves \eqref{eq:true-fold}.  The distinction is exact:
at $c=4/5$,
\[
\frac{1-c}{2}=\frac1{10},
\qquad
u_{\mathrm{fold}}(c)=\frac{392}{925}.
\]
\end{proof}

Write
\begin{align}
n(c)&:=-(1-c)c(2-c)(1+2c-c^2),\\
d(c)&:=c^3-3c+1,
\end{align}
so $u_{\mathrm{fold}}=n/d$.

\begin{proposition}[Physical fold interval]\label{prop:physical-fold}
There is a unique $c_b\in(0,1)$ satisfying
\begin{equation}\label{eq:boundary-polynomial}
p_b(c_b)=0,\qquad
p_b(c):=c^5-5c^4+6c^3-c^2+c-1.
\end{equation}
Numerically,
\[
c_b=0.613055096856276\ldots.
\]
The part of \eqref{eq:true-fold} lying in $0\leq u\leq1$ is exactly
$c_b\leq c\leq1$, with $u(c_b)=1$ and $u(1)=0$.
\end{proposition}

\begin{proof}
The equation $u_{\mathrm{fold}}=1$ is $n-d=0$, which expands to
\eqref{eq:boundary-polynomial}.  The Sturm variations of $p_b$ at $0$ and
$1$ are respectively $4$ and $3$, hence there is exactly one root in
$(0,1)$.  It is isolated by
$613/1000<c_b<614/1000$.  On $[c_b,1]$ the denominator has no zero,
$d<0$, and the inequalities $0\leq n/d\leq1$ follow from the endpoint
signs and uniqueness of the root of $n-d$.
\end{proof}

\begin{theorem}[Ordinary folds and the exceptional intersection]
\label{thm:fold-types}
Using the standard Whitney fold convention and its plane-to-plane
criterion \cite{whitney1955,golubitsky1973}, the critical points have the
following types.
\begin{enumerate}
\item Every point of $\Csym$ with $u>0$ is an ordinary Whitney fold of $F$.
\item Every point of the physical part of $\Crat$ with $c_b\leq c<1$ is
an ordinary Whitney fold, using a smooth extension across $u=1$ at the
parameter boundary.
\item At $(t,u)=(0,0)$ the two critical components meet transversely and the
ordinary-fold condition fails.
\end{enumerate}
\end{theorem}

\begin{proof}
On $\Csym$,
\[
F(0,u)=(0,1-u),\qquad F_u(0,u)=(0,-1).
\]
The kernel of $DF$ is the $t$ direction, and direct differentiation gives
$N_{tt}(0,u)=-2u$.  Hence
\[
\det\bigl(F_u,F_{tt}\bigr)=-2u\neq0
\]
for $u>0$, which is the fold criterion.

On $\Crat$, let $v=(N_u,-N_t)$ span $\ker DF$.  Since
$J_F=(-2s/Q^3)(A+uB)$, restriction to $A+uB=0$ gives
\begin{equation}\label{eq:fold-transversality}
\mathrm dJ_F(v)
=-\frac{4s^2}{Q^3B(c)}\,C_8(c),
\end{equation}
where
\[
C_8(c)=(c-1)R_7(c)
\]
and
\[
R_7(c)=3c^7-8c^6-12c^5+60c^4-69c^3+30c^2-6c-1.
\]
The Sturm variation of $R_7$ is the same at $613/1000$ and $1$, so it
has no root there; $R_7(1)=-3$.  Also $B$ has no root on this interval.
Thus \eqref{eq:fold-transversality} is nonzero for $c_b\leq c<1$.

Finally set $\delta=1-c$.  Exact substitution gives
\[
A=2\delta-3\delta^3+\delta^5,
\qquad
B=-1+3\delta^2-\delta^3.
\]
Since $\delta=t^2/2+O(t^4)$,
\begin{equation}\label{eq:degenerate-jacobian}
J_F(t,u)=2t(u-t^2)+O\!\left(|t|^5+|u||t|^3\right).
\end{equation}
The two local critical branches are therefore $t=0$ and
$u=t^2+O(t^4)$, with distinct tangent lines.  At their intersection
$DF$ has rank one but $\det(F_u,F_{tt})=0$, so the point is not an
ordinary fold.  No finer singularity label is imposed here.
\end{proof}

\begin{definition}[Observation discriminant]
\label{def:discriminant}
The observation discriminant is
\[
\DiscF:=F\bigl(\operatorname{Crit}(F)\bigr)\subset\mathbb R^2.
\]
\end{definition}

\begin{proposition}[Two critical components, two critical-value branches]
\label{prop:discriminant}
The symmetry component maps exactly to
\[
F(\Csym)=\{(N,M)=(0,m):0\leq m\leq1\}.
\]
The two signed arms of $\Crat$ have the same image because $F$ is even in
$t$.  Writing $\delta=1-c$ along that image,
\begin{align}
u_{\mathrm{fold}}&=2\delta+O(\delta^3),\\
N&=-3\delta^2+O(\delta^4),\\
M&=1-3\delta+\frac32\delta^2+O(\delta^3).
\end{align}
Consequently the rational discriminant is tangent to the symmetry segment at
$(N,M)=(0,1)$:
\begin{equation}\label{eq:discriminant-contact}
N=-\frac13(1-M)^2+O\bigl((1-M)^3\bigr).
\end{equation}
Its other physical endpoint is
\[
F(c_b,1)=
\left(
c_b^2-1,\
-\frac{(1-c_b)^2(1+c_b)}{\sqrt{c_b(2-c_b)}}
\right)
\approx(-0.6241634482,-0.2619196638).
\]
\end{proposition}

\begin{proof}
At $t=0$, substitution in $N$ and $M$ gives $(0,1-u)$.  Evenness follows
because $N$ and $M$ depend on $t$ only through $c=\cos t$ and $Q$.
The displayed expansions follow from the exact $\delta$ expressions in the
proof of \cref{thm:fold-types}.  Eliminating $\delta$ gives
\eqref{eq:discriminant-contact}; setting $u=1$ gives the final endpoint.
\end{proof}

\begin{figure}[p]
\centering
\begin{tikzpicture}
\begin{groupplot}[
  group style={group size=2 by 1,horizontal sep=1.5cm},
  width=0.46\textwidth,
  height=0.40\textwidth,
  grid=major,
  tick label style={font=\scriptsize},
  label style={font=\small},
  title style={font=\small}
]
\nextgroupplot[
  title={full critical locus},
  xlabel={$t$},ylabel={$u$},
  xmin=-0.96,xmax=0.96,ymin=-0.03,ymax=1.06
]
  \addplot[Blue,very thick] coordinates {(0,0) (0,1)};
  \addlegendentry{$\Csym$}
  \addplot[Teal,very thick] table[x=t,y=u] {\CriticalTable};
  \addplot[Teal,very thick] table[x expr=-\thisrow{t},y=u] {\CriticalTable};
  \addlegendentry{$\Crat$}
  \addplot[only marks,mark=*,mark size=2.7pt,Amber]
    coordinates {(0,0)};
\nextgroupplot[
  title={observation discriminant},
  xlabel={$N$},ylabel={$M$},
  xmin=-0.68,xmax=0.05,ymin=-0.32,ymax=1.06
]
  \addplot[Blue,very thick] coordinates {(0,0) (0,1)};
  \addlegendentry{$F(\Csym)$}
  \addplot[Coral,very thick] table[x=N,y=M] {\CriticalTable};
  \addlegendentry{$F(\Crat)$}
  \addplot[only marks,mark=*,mark size=2.7pt,Amber]
    coordinates {(0,1)};
\end{groupplot}
\end{tikzpicture}
\caption{Left: the full critical locus contains the symmetry fold $\Csym$ and
the two signed rational arms $\Crat$; they meet at the exceptional point
(amber).  Right: the symmetry fold maps to a line segment, while both rational
arms share one critical-value curve tangent to that segment at $(0,1)$.}
\label{fig:critical-discriminant}
\end{figure}

\begin{figure}[t]
\centering
\begin{tikzpicture}
\begin{axis}[
  width=0.9\textwidth,
  height=6.2cm,
  xmin=0.30,xmax=1.01,
  ymin=0,ymax=1.12,
  xlabel={$c=\cos t$},
  ylabel={$u$},
  grid=both,
  grid style={gray!20},
  legend style={at={(0.02,0.98)},anchor=north west,font=\small},
  unbounded coords=jump
]
  \path[fill=Pale] (axis cs:\cb,0) rectangle (axis cs:1,1);
  \addplot[gray!70,dashed,thick,domain=0.30:1,samples=100]
    {(1-x)/2};
  \addlegendentry{$\LN:\ u=(1-c)/2$}
  \addplot[Coral,thick,domain=0.3476:1,samples=300,
           restrict y to domain=0:1.12]
    {ufold(x)};
  \addlegendentry{algebraic continuation of $\Crat$}
  \addplot[Teal,very thick,domain=\cb:1,samples=180]
    {ufold(x)};
  \addlegendentry{physical part of $\Crat$}
  \addplot[Blue,densely dotted] coordinates {(\cb,0) (\cb,1.12)};
  \node[Blue,anchor=south west,font=\small]
    at (axis cs:\cb,0.03) {$c_b$};
  \addplot[only marks,mark=*,mark size=2.5pt,Amber]
    coordinates {(0.38196601125,0.30901699437)};
  \node[Amber,anchor=south west,font=\scriptsize]
    at (axis cs:0.38196601125,0.30901699437) {lateral zero};
\end{axis}
\end{tikzpicture}
\caption{The component-zero curve $\LN$ and the rational critical component
$\Crat$ are different.  The thick teal segment is the complete portion of
$\Crat$ inside the parameter strip.  The symmetry component $\Csym$, visible
in $(t,u)$ coordinates in \cref{fig:critical-discriminant}, collapses to the
boundary line $c=1$ in this plot.}
\label{fig:locus-audit}
\end{figure}

\clearpage
\section{Spherical radial image of the rational component}

\begin{definition}[Spherical radial fold map]\label{def:stokes-map}
For $c\in[c_b,1]$, put $u=u_{\mathrm{fold}}(c)$ and define
\begin{align}
X(c)&=c+2u(1-c),\\
Y_\pm(c)&=\pm(1-u)\sqrt{1-c^2},\\
Z(c)&=u\sqrt{c(2-c)},\\
R(c)&=\sqrt{X(c)^2+Y_\pm(c)^2+Z(c)^2}.
\end{align}
The two normalized branches are
\begin{equation}\label{eq:corrected-stokes}
\Phi_\pm(c):=\frac{1}{R(c)}
\bigl(X(c),Y_\pm(c),Z(c)\bigr)\in S^2.
\end{equation}
\end{definition}

The sign distinguishes $t>0$ and $t<0$; $R$ is independent of that sign.
We use the conventional coordinate labels $(S_1,S_2,S_3)$ when plotting this
unit vector on the Stokes--Poincar\'e sphere
\cite{stokes1852,poincare1892}, but no Jones-vector identity is assumed.

\begin{remark}[Three maps, no automatic bridge]\label{rem:three-maps}
The observation discriminant $\DiscF\subset\mathbb R^2$, the radial map
$\Phi=S/\lVert S\rVert$, and the Gauss map
\[
\nu=\frac{S_t\times S_u}{\lVert S_t\times S_u\rVert}
\]
are distinct constructions.  In particular, translating the surface changes
$\Phi$ but leaves $F$, $\DiscF$ and $\nu$ unchanged.  Therefore neither a
discriminant interpretation nor a polarization interpretation of $\Phi$
follows from normalization alone.
\end{remark}

\begin{theorem}[Endpoints]\label{thm:endpoints}
The two branches have common endpoints
\begin{align}
\Phi_\pm(1)&=(1,0,0),\\
\Phi_\pm(c_b)&=
\left(
\sqrt{\frac{2-c_b}{2}},
0,
\sqrt{\frac{c_b}{2}}
\right)\\
&\approx(0.832749933396492,0,0.553649300937099).
\end{align}
\end{theorem}

\begin{proof}
At $c=1$, Proposition~\ref{prop:physical-fold} gives $u=0$, and
\eqref{eq:surface} equals $(1,0,0)$.  At $c=c_b$, $u=1$, so
$S=(2-c_b,0,\sqrt{c_b(2-c_b)})$.  Its squared norm is
$2(2-c_b)$, which gives the displayed unit vector.
\end{proof}

\begin{theorem}[Regular radial starting point]\label{thm:regular-start}
The spherical radial fold map is regular at $t=0$.  With $c=\cos t$ and the sign of
$Y$ inherited from $t$,
\begin{equation}\label{eq:start-expansion}
\Phi(t)=
\left(
1-\frac{t^2}{2}+O(t^4),
\ t+O(t^3),
\ t^2+O(t^4)
\right).
\end{equation}
Consequently
\[
\Phi'(0)=(0,1,0)\neq0.
\]
\end{theorem}

\begin{proof}
Expansion of \eqref{eq:true-fold} at $c=\cos t$ gives
$u_{\mathrm{fold}}(\cos t)=t^2+O(t^4)$.  Substitution in
\eqref{eq:surface}, followed by normalization, gives
\eqref{eq:start-expansion}.
\end{proof}

\begin{remark}
The former argument tested whether a one-dimensional curve surjects onto the
two-dimensional tangent plane of $S^2$.  That is not the immersion
criterion.  Rank one is maximal for a curve, and
\cref{thm:regular-start} supplies the required nonzero derivative.
\end{remark}

\begin{figure}[p]
\centering
\begin{tikzpicture}
\begin{axis}[
  width=0.78\textwidth,
  height=8.2cm,
  view={-52}{22},
  xlabel={$S_1$},ylabel={$S_2$},zlabel={$S_3$},
  xmin=0.80,xmax=1.01,
  ymin=-0.46,ymax=0.46,
  zmin=0,zmax=0.59,
  grid=major,
  legend style={at={(0.02,0.98)},anchor=north west,font=\small}
]
  \addplot3[Teal,very thick]
    table[x=Sone,y=Stwo,z=Sthree] {\StokesTable};
  \addlegendentry{$t>0$}
  \addplot3[Coral,very thick]
    table[x=Sone,y expr=-\thisrow{Stwo},z=Sthree] {\StokesTable};
  \addlegendentry{$t<0$}
  \addplot3[only marks,mark=*,mark size=2.8pt,Blue]
    coordinates {(1,0,0)};
  \addplot3[only marks,mark=*,mark size=2.8pt,Amber]
    coordinates {(0.832749933396492,0,0.553649300937099)};
\end{axis}
\end{tikzpicture}
\caption{Spherical radial image of the physical part of $\Crat$.  The two
branches are mirror images across $S_2=0$ and share both the regular radial
starting point (blue) and the parameter-boundary endpoint (amber).  This is
not the planar discriminant shown in \cref{fig:critical-discriminant}.}
\label{fig:stokes-3d}
\end{figure}

\begin{figure}[p]
\centering
\begin{tikzpicture}
\begin{groupplot}[
  group style={group size=3 by 1,horizontal sep=1.1cm},
  width=0.31\textwidth,
  height=0.31\textwidth,
  grid=major,
  tick label style={font=\scriptsize},
  label style={font=\small},
  title style={font=\small}
]
\nextgroupplot[
  title={top view},
  xlabel={$S_1$},ylabel={$S_2$},
  xmin=0.80,xmax=1.01,ymin=-0.46,ymax=0.46
]
  \addplot[Teal,very thick]
    table[x=Sone,y=Stwo] {\StokesTable};
  \addplot[Coral,very thick]
    table[x=Sone,y expr=-\thisrow{Stwo}] {\StokesTable};
  \addplot[only marks,mark=*,Blue] coordinates {(1,0)};
  \addplot[only marks,mark=*,Amber]
    coordinates {(0.832749933396492,0)};
\nextgroupplot[
  title={front view},
  xlabel={$S_1$},ylabel={$S_3$},
  xmin=0.80,xmax=1.01,ymin=0,ymax=0.59
]
  \addplot[Teal,very thick]
    table[x=Sone,y=Sthree] {\StokesTable};
  \addplot[only marks,mark=*,Blue] coordinates {(1,0)};
  \addplot[only marks,mark=*,Amber]
    coordinates {(0.832749933396492,0.553649300937099)};
\nextgroupplot[
  title={side view},
  xlabel={$S_2$},ylabel={$S_3$},
  xmin=-0.46,xmax=0.46,ymin=0,ymax=0.59
]
  \addplot[Teal,very thick]
    table[x=Stwo,y=Sthree] {\StokesTable};
  \addplot[Coral,very thick]
    table[x expr=-\thisrow{Stwo},y=Sthree] {\StokesTable};
  \addplot[only marks,mark=*,Blue] coordinates {(0,0)};
  \addplot[only marks,mark=*,Amber]
    coordinates {(0,0.553649300937099)};
\end{groupplot}
\end{tikzpicture}
\caption{Three orthogonal views of the spherical radial image.  These projections
make the mirror symmetry and common endpoints visible without relying on a
particular three-dimensional camera angle.}
\label{fig:orthogonal-views}
\end{figure}

\clearpage
\section{Boundary endpoint and third spherical component}

Let $\mathcal C(c):=S_3(\Phi_+(c))=Z(c)/R(c)$.

\begin{theorem}[Boundary maximum]\label{thm:boundary-maximum}
On the physical fold interval $[c_b,1]$, $\mathcal C(c)$ decreases strictly
with $c$.  Equivalently it increases from the regular start toward the
parameter boundary.  Hence
\begin{equation}
\max_{c_b\leq c\leq1}\mathcal C(c)
=\mathcal C(c_b)
=\sqrt{\frac{c_b}{2}}
\approx0.553649300937099.
\end{equation}
There is no interior critical point of the third spherical component on the
physical part of $\Crat$.
\end{theorem}

\begin{proof}
Writing $\mathcal C^2=A/R_{\mathrm{num}}$ with exact integer polynomials,
direct differentiation gives
\begin{equation}
\frac{d}{dc}\mathcal C(c)^2
=\frac{n(c)Q_{17}(c)}{R_{\mathrm{num}}(c)^2}.
\end{equation}
The polynomial $Q_{17}$ is given in
\cref{app:certificate}.  Its Sturm variation is $9$ at both $613/1000$ and
$1$, so it has no root on that interval; $Q_{17}(4/5)>0$.  Since
$n(c)<0$ for $c_b<c<1$, the derivative is negative.  The endpoint value is
Theorem~\ref{thm:endpoints}.
\end{proof}

\begin{figure}[t]
\centering
\begin{tikzpicture}
\begin{axis}[
  width=0.82\textwidth,
  height=6cm,
  xmin=0.60,xmax=1.01,
  ymin=0,ymax=0.59,
  xlabel={$c=\cos t$},
  ylabel={$\mathcal C(c)=S_3$},
  grid=both,
  grid style={gray!20}
]
  \addplot[Teal,very thick]
    table[x=c,y=Sthree] {\StokesTable};
  \addplot[only marks,mark=*,mark size=2.8pt,Amber]
    coordinates {(\cb,0.553649300937099)};
  \addplot[Blue,densely dotted]
    coordinates {(\cabs,0) (\cabs,0.59)};
  \node[Blue,anchor=south west,font=\small]
    at (axis cs:\cabs,0.03) {$c^{**}$};
  \node[Amber,anchor=north west,font=\small]
    at (axis cs:\cb,0.553649300937099) {boundary maximum};
\end{axis}
\end{tikzpicture}
\caption{The third component of the spherical radial image.  The former
critical value $c_{\mathrm{chiral}}\approx0.3103$ lies outside the physical
interval and is therefore absent.  The dotted line marks the
rational-critical-branch--skeleton intersection
$c^{**}\approx0.673741$.  A circular-polarization interpretation is not
assumed.}
\label{fig:chirality}
\end{figure}

\begin{remark}[What this boundary is not]
The point $u=1$ is a boundary point of the two-dimensional ruled-surface
parameter strip.  The construction supplies neither a spinor bundle on
$S^3$, a boundary Dirac operator, nor a self-adjoint boundary domain.
Consequently it cannot by itself define an APS spectral projector.  Any such
identification would be an additional bridge.
\end{remark}

\clearpage
\section{Disposition of earlier claims}

\begin{center}
\begin{tabular}{@{}p{0.35\textwidth}p{0.55\textwidth}@{}}
\toprule
Earlier statement & Status in v5\\
\midrule
$u_*=(1-c)/2$ is the Whitney fold
& False.  It is the component-zero curve $\LN$.\\
The rational branch is the full locus $\{J_F=0\}$
& False.  The symmetry component $\Csym$ was omitted.\\
$\Phi=S/\lVert S\rVert$ is the observation discriminant
& False without a bridge; the discriminant is $F(\operatorname{Crit}(F))$.\\
The normalized curve is singular at $t=0$
& False.  Its derivative is $(0,1,0)$.\\
Endpoint $(2/\sqrt5,\pm1/\sqrt5,0)$
& Belongs to the auxiliary $\LN$ curve, not the true fold.\\
Irreducible chirality quintic
& Describes an interior maximum of that auxiliary curve; it is outside the
  physical true-fold interval.\\
$\sqrt5$ endpoint thread
& An identity for the auxiliary curve only; removed from the radial-fold
  theorem.\\
\bottomrule
\end{tabular}
\end{center}

The auxiliary $\LN$ construction can still be studied under its correct
name, but none of its constants may be transferred to the observation
discriminant or radial fold image without an additional theorem.

\section{Conclusion}

The observation map has a two-component critical locus.  The symmetry segment
$\Csym$ consists of ordinary folds for $u>0$; the physical rational component
$\Crat$ consists of ordinary folds away from its common initial point.  At
$(0,0)$ the components meet, the Jacobian has the local factor
$2t(u-t^2)$, and ordinary-fold regularity fails.  The discriminant is the
union of a line segment and a rational critical-value curve with quadratic
tangency at $(0,1)$.

The separately defined spherical radial image retains the exact physical
interval, endpoints, nonzero starting tangent and monotone third component
proved in v4.  It is origin-dependent and is not identified with the
observation discriminant or the Gauss map.  Thus no Jones, spinorial, quantum
or APS conclusion follows from the present construction.  A future Hopf lift
must first choose and justify which sphere map is being lifted.

There remains one local singularity-theory problem.  Put $x=1-M$ and $y=t$.
Since $x_u(0,0)=1$, $(x,y)$ are valid local source coordinates.  With weights
$\operatorname{wt}(y)=1$ and $\operatorname{wt}(u)=2$, direct expansion gives
\[
x=u+\frac12y^2+O_{\mathrm w}(4),\qquad
N=-uy^2+\frac14y^4+O_{\mathrm w}(6),
\]
and hence, after eliminating $u$,
\[
(x,N)=\left(x,-xy^2+\frac34y^4+O_{\mathrm w}(6)\right).
\]
Its displayed four-jet is equivalent by real rescalings to the specified
jet $(x,xy^2+y^4)$ occurring in the plane-to-plane recognition problem
\cite{kabata2016}.  This does not classify the full germ: the exact map is
even under $t\mapsto-t$, and finite determinacy and the higher-jet recognition
conditions remain to be checked.  The next well-posed problem is therefore an
$\mathcal A$-classification, or a symmetry-preserving $\mathbb Z_2$
classification, followed only then by a versal symmetry-breaking unfolding.

\appendix
\section{Exact verification certificate}\label{app:certificate}

The companion file \texttt{verify\_stokes\_caustic\_v5.py} uses only the
Python standard library.  It checks the exact Jacobian factors, the local
$\delta=1-c$ expansions, the weighted map-germ jet, Sturm root counts, fold
transversality, discriminant tangency, the physical interval, unit-sphere
normalization, endpoint values, the nonzero radial starting derivative and
third-component monotonicity.

For reproducibility, the derivative factor in
\cref{thm:boundary-maximum} is
\[
Q_{17}(c)=\sum_{j=0}^{17}q_jc^j,
\]
with coefficients in increasing order
\[
\begin{split}
(q_0,\ldots,q_{17})={}&
(0,-12,14,64,-346,856,-810,798,-4868,\\
&14160,-21630,20432,-12800,5462,-1588,306,-36,2).
\end{split}
\]
The exact Sturm counts used in the paper are
\[
\begin{array}{c|cc|c}
\text{polynomial}&V(\text{left})&V(\text{right})&
\text{roots in interval}\\
\hline
p_b,\ (0,1)&4&3&1\\
R_7,\ (613/1000,1)&3&3&0\\
Q_{17},\ (613/1000,1)&9&9&0.
\end{array}
\]

\begin{thebibliography}{9}
\bibitem{hayami2026surface}
A.~Hayami (J.-Y.~Huang),
``The orthogonal-circle ruled surface: Immersion theory, a topological
dipole, and projection singularities,'' companion preprint, version of
June 2026.

\bibitem{whitney1955}
H.~Whitney,
``On singularities of mappings of Euclidean spaces. I. Mappings of the plane
into the plane,''
\emph{Annals of Mathematics}, vol.~62, no.~3, pp.~374--410, 1955,
doi:10.2307/1970070.

\bibitem{golubitsky1973}
M.~Golubitsky and V.~Guillemin,
\emph{Stable Mappings and Their Singularities}, Graduate Texts in
Mathematics, vol.~14. Springer, 1973,
doi:10.1007/978-1-4615-7904-5.

\bibitem{kabata2016}
Y.~Kabata,
``Recognition of plane-to-plane map-germs,''
\emph{Topology and its Applications}, vol.~202, pp.~216--238, 2016,
doi:10.1016/j.topol.2016.01.011.

\bibitem{stokes1852}
G.~G.~Stokes,
``On the composition and resolution of streams of polarized light from
different sources,''
\emph{Transactions of the Cambridge Philosophical Society}, vol.~9,
pp.~399--416, 1852.

\bibitem{poincare1892}
H.~Poincar\'e,
\emph{Th\'eorie Math\'ematique de la Lumi\`ere}, vol.~2.
Gauthier--Villars, 1892.
\end{thebibliography}

\end{document}
